% ==============================================================================
% PHỤ LỤC: TRUY NGUỒN SỐ LIỆU VÀ TÁI LẬP
% ==============================================================================
\needspace{0.6\textheight}
\section{Truy nguồn số liệu và cách tái lập}
\label{sec:trace}

Mọi con số trong báo cáo được tính bằng script của nhóm từ dữ liệu lưu trong repo (mục~\ref{sec:storage}); không có số liệu nào được nhập thủ công. Bảng~\ref{tab:scripts} cho biết mỗi nhóm số liệu được tạo ra ở bước nào.

\begin{table}[H]
\centering
\caption{Nguồn gốc của các nhóm số liệu trong báo cáo}
\label{tab:scripts}
\begin{tabularx}{\textwidth}{>{\raggedright\arraybackslash}p{0.36\linewidth} L}
\toprule
\textbf{Số liệu} & \textbf{Được tạo ra bởi} \\
\midrule
Số ảnh Mapillary, số đường OSM, tỉ lệ gắn đường, số ảnh đêm bị loại (mục~\ref{sec:pipeline}) & Các bước thu thập; kết quả của từng bước được ghi vào nhật ký thu thập và bảng trạng thái thu thập \\
So sánh ảnh bên trái và bên phải, tham số cắt ảnh & Hai phép đo thăm dò về độ phủ vỉa hè theo phía chụp và theo góc cắt \\
Đồng thuận, phân xử, nguồn nhãn (mục~\ref{sec:qa}) & Bước đo đồng thuận và dựng nhãn cuối từ phiếu gán và bảng phân xử \\
Bộ chọn ảnh, chọn đường, chọn 2.000 ảnh bổ sung (mục~\ref{sec:enrich}) & Phép so sánh các bộ chọn ảnh (để riêng từng tuyến đường), lượt quét 150 đường và bước chọn ảnh bổ sung \\
Kết quả gán nhãn ảnh bổ sung, kiểm tra chéo & Bước tổng hợp nhãn ảnh bổ sung từ phiếu gán và bảng kiểm tra chéo \\
Bảng~\ref{tab:overview}, \ref{tab:label-source}, \ref{tab:inventory}--\ref{tab:assoc}, \ref{tab:checks}; các hình về nhãn và metadata & Script EDA dữ liệu dạng bảng: đếm nhãn, khoảng tin cậy, thống kê mô tả, liên hệ với nhãn và 71 phép kiểm tra nhất quán; script này chỉ đọc bảng nhãn và manifest \\
Bảng~\ref{tab:image}; các hình về chất lượng ảnh và ảnh ví dụ & Script EDA ảnh: đo từng ảnh cắt (độ sáng, độ tương phản, tỉ lệ điểm ảnh tối và cháy sáng, độ nét, độ sặc sỡ màu) \\
Bảng~\ref{tab:spatial}, \ref{tab:domain}; các hình về trùng lặp, vị trí và t-SNE & Script EDA trùng lặp và rò rỉ: mã băm, dHash, embedding DINOv2, khoảng cách địa lý giữa các tập, bộ phân loại miền \\
\bottomrule
\end{tabularx}
\end{table}

\textbf{Tái lập kết quả.} Ảnh không được lưu trong git. Có hai cách lấy ảnh: dùng bản ảnh cắt trên server gán nhãn, hoặc dựng lại từ manifest (mã ảnh gốc và góc quay) bằng script dựng lại ảnh của nhóm, với điều kiện có token Mapillary. Bản dựng lại khác bản trên server do nén JPEG lại, nên cần ghi rõ nguồn ảnh đã dùng; 671 ảnh val/test trong phần EDA ảnh của báo cáo này là bản dựng lại. Tiếp theo, cài đặt các thư viện trong danh sách phụ thuộc của repo, ghi lại phiên bản Python và thư viện thực sự sử dụng, rồi chạy lần lượt ba script EDA (dạng bảng, ảnh, trùng lặp). Script dạng bảng chỉ cần các bảng dữ liệu; hai script còn lại cần kho ảnh cắt, và script trùng lặp nên được chạy trên GPU. Báo cáo được biên dịch bằng XeLaTeX và BibTeX.

\textbf{Phân công và việc sử dụng công cụ AI.} Ba script EDA do ba tác tử AI (Claude Sonnet) viết theo yêu cầu của một tác tử điều phối (Claude). Tác tử điều phối đã tính lại độc lập các số liệu chính (đếm nhãn, khoảng Wilson, thống kê metadata, phép đo ảnh trên 300 ảnh mẫu, mã băm MD5, khoảng cách giữa các tập, độ tương đồng embedding) và xem các bảng ảnh kiểm tra trước khi đưa kết quả vào báo cáo. Các kết quả kiểm tra trực quan trong báo cáo là quan sát của AI trên các bảng ảnh tổng hợp, không phải nhãn và không phải phép đo. Nhãn của bộ dữ liệu do người gán và người phân xử quyết định: Claude chỉ đề xuất nhãn cho 303 ảnh bất đồng của bộ chính (do Khải duyệt) và 150 ảnh pilot (do Phúc duyệt); 7 nhãn cuối của bộ chính là các đề xuất khác cả hai phiếu và được người phân xử giữ lại (mục~\ref{sec:qa}).

\textbf{Ảnh minh họa trong báo cáo.} Hình~\ref{fig:examples} và~\ref{fig:examples-extra} chỉ dùng ảnh train. Các bảng ảnh kiểm tra nội bộ có chứa ảnh val/test nên không được đưa vào báo cáo và không được đưa lên git. Nguồn ảnh: \cc.
